\documentclass[10pt,a4paper]{amsart}
\usepackage{amsmath,amssymb,amsthm,graphicx}
% Retained preamble for the Stacks Project chapter snapshot in this directory.
% The chapter's own preamble.tex is a copy of the shared Stacks Project preamble; the block
% below is sliced MECHANICALLY, line for line, from it. Nothing was added, removed or
% redefined; the wrapper class, the xr-hyper/hyperref wiring and the externaldocument list
% are omitted because the assembler supplies the document class.
%
% lines taken from preamble.tex: 3, 5, 9, 11, 12, 156, 157, 158, ... (51 lines total)
%
% package coverage restored from the chapter's own preamble (2026-10-13): xy, xy-2cell

\usepackage[all]{xy}
\xyoption{2cell}
\UseAllTwocells

\newenvironment{reference}{\comment}{\endcomment}
\newenvironment{slogan}{\comment}{\endcomment}
\newenvironment{history}{\comment}{\endcomment}
\theoremstyle{plain}
\newtheorem{theorem}[subsection]{Theorem}
\newtheorem{proposition}[subsection]{Proposition}
\newtheorem{lemma}[subsection]{Lemma}
\theoremstyle{definition}
\newtheorem{definition}[subsection]{Definition}
\newtheorem{example}[subsection]{Example}
\newtheorem{exercise}[subsection]{Exercise}
\newtheorem{situation}[subsection]{Situation}
\theoremstyle{remark}
\newtheorem{remark}[subsection]{Remark}
\newtheorem{remarks}[subsection]{Remarks}
\numberwithin{equation}{subsection}
\def\lim{\mathop{\mathrm{lim}}\nolimits}
\def\colim{\mathop{\mathrm{colim}}\nolimits}
\def\Spec{\mathop{\mathrm{Spec}}}
\def\Hom{\mathop{\mathrm{Hom}}\nolimits}
\def\Ext{\mathop{\mathrm{Ext}}\nolimits}
\def\SheafHom{\mathop{\mathcal{H}\!\mathit{om}}\nolimits}
\def\SheafExt{\mathop{\mathcal{E}\!\mathit{xt}}\nolimits}
\def\Sch{\mathit{Sch}}
\def\Mor{\mathop{\mathrm{Mor}}\nolimits}
\def\Ob{\mathop{\mathrm{Ob}}\nolimits}
\def\Sh{\mathop{\mathit{Sh}}\nolimits}
\def\NL{\mathop{N\!L}\nolimits}
\def\CH{\mathop{\mathrm{CH}}\nolimits}
\def\proetale{{pro\text{-}\acute{e}tale}}
\def\etale{{\acute{e}tale}}
\def\QCoh{\mathit{QCoh}}
\def\Ker{\mathop{\mathrm{Ker}}}
\def\Im{\mathop{\mathrm{Im}}}
\def\Coker{\mathop{\mathrm{Coker}}}
\def\Coim{\mathop{\mathrm{Coim}}}
\DeclareMathSymbol{\boxtimes}{\mathbin}{AMSa}{"02}
\def\QCohstack{\mathcal{QC}\!\mathit{oh}}
\def\Cohstack{\mathcal{C}\!\mathit{oh}}
\def\Spacesstack{\mathcal{S}\!\mathit{paces}}
\def\Quotfunctor{\mathrm{Quot}}
\def\Hilbfunctor{\mathrm{Hilb}}
\def\Curvesstack{\mathcal{C}\!\mathit{urves}}
\def\Polarizedstack{\mathcal{P}\!\mathit{olarized}}
\def\Complexesstack{\mathcal{C}\!\mathit{omplexes}}
\def\Pic{\mathop{\mathrm{Pic}}\nolimits}
\def\Picardstack{\mathcal{P}\!\mathit{ic}}
\def\Picardfunctor{\mathrm{Pic}}
\def\Deformationcategory{\mathcal{D}\!\mathit{ef}}
% lane-s11: discard bibliographic \begin{reference} blocks entirely (incl. any \cite inside).
\newsavebox{\limrefbox}
\ifcsname reference\endcsname\expandafter\let\csname reference\endcsname\undefined\fi\ifcsname endreference\endcsname\expandafter\let\csname endreference\endcsname\undefined\fi
\def\citation#1{\relax}
\newenvironment{reference}{\setbox\limrefbox=\vbox\bgroup\let\cite\relax}{\egroup\ignorespaces}
\usepackage{cleveref}
\crefname{originalreference}{original reference}{original references}
\Crefname{originalreference}{Original reference}{Original references}
\makeatletter
\newlabel{situation-descent-property}{{081D}{source}{}{sourcepointer-1}{}}
\newlabel{situation-descent-property@cref}{{[originalreference][0][]081D}{[1][0][]source}{}{}{}}
\makeatother
\crefname{originalreference}{original reference}{original references}
\Crefname{originalreference}{Original reference}{Original references}
\makeatletter
\newlabel{algebra-lemma-flat-finite-presentation-limit-flat}{{02JO}{source}{}{sourcepointer-2}{}}
\newlabel{algebra-lemma-flat-finite-presentation-limit-flat@cref}{{[originalreference][0][]02JO}{[1][0][]source}{}{}{}}
\makeatother
\title{Descent of flat finite presentation morphisms along limits}
\author{The Stacks Project Authors}
\date{stacks-project:limits.tex:local-snapshot-20261012; GNU FDL 1.2 or later}
\begin{document}
\maketitle
\noindent\textit{Editorial note.} No omissions in the counted span: the reduction is worked; the algebra core is a precise citation of a proved lemma.\par\medskip
\noindent\textit{Allowed standard prerequisites (editorial).}\begin{itemize}
\item Limit-flat finite presentation. Algebra, Lemma flat-finite-presentation-limit-flat
\end{itemize}\medskip
\ifdefined\hypertarget\hypertarget{sourcepointer-1}{}\fi
\noindent\textit{Source reference.} Original reference 081D: sources/stacks-chow/tags.txt, source lines 10213--10213. Cross-chapter standard prerequisite cited by the author inside the counted span; it is used as a pointer and its own proof is not part of this item.\par\medskip
\ifdefined\hypertarget\hypertarget{sourcepointer-2}{}\fi
\noindent\textit{Source reference.} Original reference 02JO: sources/stacks-chow/tags.txt, source lines 3343--3343. Cross-chapter standard prerequisite cited by the author inside the counted span; it is used as a pointer and its own proof is not part of this item.\par\medskip
\begin{lemma}
\label{lemma-descend-flat-finite-presentation}
Notation and assumptions as in Situation \ref{situation-descent-property}.
If
\begin{enumerate}
\item $f$ is flat,
\item $f_0$ is locally of finite presentation,
\end{enumerate}
then $f_i$ is flat for some $i \geq 0$.
\end{lemma}

\begin{proof}
Choose a finite affine open covering
$Y_0 = \bigcup_{j = 1, \ldots, m} Y_{j, 0}$
such that each $Y_{j, 0}$ maps into an affine open
$S_{j, 0} \subset S_0$. For each $j$ let
$f_0^{-1}Y_{j, 0} = \bigcup_{k = 1, \ldots, n_j} X_{k, 0}$ be a
finite affine open covering. Since the property of being flat is
local we see that it suffices to prove the lemma for the morphisms
of affines $X_{k, i} \to Y_{j, i} \to S_{j, i}$ which are the base
changes of $X_{k, 0} \to Y_{j, 0} \to S_{j, 0}$ to $S_i$.
Thus we reduce to the case that $X_0, Y_0, S_0$ are affine

\medskip\noindent
In the affine case we reduce to the following algebra result.
Suppose that $R = \colim_{i \in I} R_i$. For some $0 \in I$
suppose given an $R_0$-algebra map $A_i \to B_i$ of finite presentation.
If $R \otimes_{R_0} A_0 \to R \otimes_{R_0} B_0$ is flat, then
for some $i \geq 0$ the map
$R_i \otimes_{R_0} A_0 \to R_i \otimes_{R_0} B_0$ is flat.
This follows from Algebra,
Lemma \ref{algebra-lemma-flat-finite-presentation-limit-flat} part (3).
\end{proof}

\end{document}
